\chapter{网格自适应性}

\section{网格自适应加密放粗}

网格自适应可以分为局部加密放粗和全局加密放粗，局部和全局的区别关键在于对初始网格的处理，局部网格剖分是对初始网格的修改，而全局网格剖分会重新生成新的网格。
局部加密放粗可以分为按顶点加密放粗和按单元加密放粗，全局加密放粗可以采用quadtree-octree（四叉树/八叉树）、advancing-front（前沿推进）和Delaunay类型的。
图\ref{fig:2-1}中用两个简单的例子说明按顶点加密放粗和按单元加密放粗，左上分图里红点为顶点，围绕该顶点有两条边，对两条边取中点得到两个新点，并和原有顶点相连，进行了细分，
左下分图里两个红点代表两个单元，将两个单元相邻边取中点，和边对面的顶点相连，分别将两个单元进行了细分得到了四个单元，并且保持协调性。全局加密放粗计算量大于局部，此外还有
hierarchic methods（继承方法）、多重网格、非协调和重叠方法，还需要考虑各向异性。

\begin{figure}[!htbp]
  \centering
  \includegraphics[height=5cm]{fig/2/1.png}
  \caption{按定点加密放粗和按单元加密放粗}
  \label{fig:2-1}
\end{figure}


\section{层网格自适应加密放粗}
